\section{Evaluation}
\label{sec:eval}

We ask five questions.
\begin{itemize}
  \item[\rqref{1}] Does \tool (the invariant-based Protocol B engine) detect
    the motivating CVE on real OpenSSL source, distinguish the fix from the
    bug, and generalize to a fix it was not derived from?
  \item[\rqref{2}] Do general-purpose open-source static-analysis tools
    already catch the bugs a source-pattern baseline catches, on the same
    three CVEs?
  \item[\rqref{3}] Does a rule written from one CVE flag a site that
    OpenSSL fixed only later, without the rule being changed after the
    fact?
  \item[\rqref{4}] What does \tool cost in time?
  \item[\rqref{5}] Does \tool's rule-as-data design extend to a second,
    differently-shaped invariant (Protocol A, CVE-2023-0215), and does that
    second rule generalize the same way Protocol B's did?
\end{itemize}

\subsection{Subjects}
\label{sec:subjects}
We select three historical OpenSSL CVEs and one non-CVE fix, each a pinned
vulnerable/fixed commit pair with a publicly documented root cause (so the
label is not our judgment), source at both revisions in the same file, and
at least one case not used to derive the rule (\rqref{3}).
CVE-2022-4450 (\texttt{PEM\_read\_bio\_ex}, fix \texttt{63bcf189be}) is the
source of Protocol~B. CVE-2023-0215 (\texttt{BIO\_new\_NDEF}, fix
\texttt{8818064ce3}) is the second baseline CVE (\rqref{2}) and the source
of Protocol~A (\rqref{5}). CVE-2024-4741 (\texttt{SSL\_free\_buffers}, fix
\texttt{b3f0eb0a29}) is in the baseline comparison only: it needs a state
guard neither rule can express, and we report that scope rather than drop
the case. The \texttt{d2217c88df} KDF fix (not a CVE, a different author, a
year before CVE-2022-4450) is the out-of-sample case.

\subsection{Baselines}
\label{sec:baselines}
The \emph{Part~1 pattern baseline} is \texttt{typestate\_check.py}, the
project's own string/regex patch matcher, included as the simplest thing
that could work. \emph{Clang Static Analyzer~18} ran via
\texttt{scan-build-18} on each revision's affected translation unit with its
generated build flags. \emph{Coccinelle~1.1.1} ran five pre-existing
semantic patches (not written for these CVEs) for release/use, repeated
release, release-after-transfer, and allocate-then-return.
\emph{Cppcheck~2.13.0} ran with \texttt{--enable=warning,style,performance,
portability --inconclusive} (C11/unix64). None was tuned toward these
CVEs. Because Protocol~B has no rule for two of the three CVEs, this
compares general tools against the baseline, not against \tool. For each
tool and revision we manually read whether the output distinguishes
vulnerable from fixed.

\subsection{Setup}
\label{sec:setup}
\tool (Protocol~B) ran on Python~3.13, \texttt{tree-sitter}~0.26.0,
\texttt{tree-sitter-c}~0.24.2 (Linux, WSL2). The \rqref{2} baselines ran on
a separate machine (Python~3.12.3, Clang~18.1.3, Coccinelle~1.1.1 on
OCaml~4.14.1, Cppcheck~2.13.0); executable hashes and frozen rule files are
recorded with the data, and no measurement mixes the two environments. For
\rqref{4}, a second author reran \texttt{run\_scan.sh} on a third machine
(Python~3.11.6, same \texttt{tree-sitter} versions) under
\texttt{/usr/bin/time -v}. Protocol~A's runs (\rqref{5}) were reproduced on
a fourth environment (\texttt{tree-sitter}~0.26.1) with identical counts.

\subsection{\rqref{1}: Does \tool detect and generalize the motivating CVE?}
\label{sec:rq1}

\begin{table}[t]
  \centering
  \caption{\tool (Protocol B) on six pinned OpenSSL snapshots
    (\texttt{crypto/}, \texttt{ssl/}, \texttt{providers/}). A
    \emph{release site} is a call that frees \texttt{*P} for a
    pointer-to-pointer parameter \texttt{P}. The two fixes each remove
    exactly their own violations and add none.}
  \label{tab:rq1-wholetree}
  \begin{tabular}{lrrrr}
    \toprule
    Snapshot & C files & Conforming & Violation & Suppressed \\
    \midrule
    CVE-2022-4450 pre-fix (\texttt{8e257b86e5})  & 1001 & 99  & 7  & 2 \\
    CVE-2022-4450 fix (\texttt{63bcf189be})      & 1001 & 101 & 5  & 2 \\
    KDF fix parent (\texttt{29a27cb2c5})         & 1005 & 93  & 12 & 2 \\
    KDF fix (\texttt{d2217c88df}, not a CVE)     & 1005 & 97  & 8  & 2 \\
    OpenSSL 3.0.0 (\texttt{89cd17a031})          & 999  & 92  & 10 & 2 \\
    master (\texttt{949fe8b5f4}, 2026-10-30)     & 1155 & 132 & 10 & 2 \\
    \bottomrule
  \end{tabular}
\end{table}

Table~\ref{tab:rq1-wholetree} shows \tool on the motivating function and on
the whole tree. At the vulnerable revision, \tool flags
\texttt{PEM\_read\_bio\_ex}'s \texttt{*header} and \texttt{*data} as
violations (witness lines 991--998); at the fix, both sites are
conforming. The same fix removes exactly these two violations from the
whole-tree count (99~$\to$~101 conforming, 7~$\to$~5 violation) and adds
none.

\begin{observation}
\label{obs:rq1-conformance}
\tool is not limited to the function it was derived from: the unrelated,
non-CVE \texttt{d2217c88df} KDF fix removes exactly the four violations
its own commit changes (\texttt{pbkdf1}, \texttt{pbkdf2}, \texttt{pkcs12kdf},
\texttt{scrypt}, each \texttt{*buffer}) and adds none, with the rule
unchanged from the CVE-2022-4450 derivation.
\end{observation}

\begin{observation}
\label{obs:rq1-mastertriage}
On current master, \tool flags 10 sites; manual triage against each site's
callers found 2 potential recurrences whose failure branch is not
reachable in normal builds, 8 convention deviations with no caller impact,
and 0 confirmed new vulnerabilities. (Two of the eight were initially
labelled "unclear" pending a caller read; Section~\ref{sec:threats}
reports that re-check.) We report this plainly: the contribution here is
a bounded rule that reproduces two real fixes without retuning, not a new
bug.
\end{observation}

\subsection{\rqref{2}: Do general-purpose tools already catch these bugs?}
\label{sec:rq2}

\begin{table}[t]
  \centering
  \caption{Detection on the vulnerable/fixed pair for each CVE: whether
    the tool's output distinguishes the two revisions for the
    CVE's own mechanism. \cmark{} = distinguishes vulnerable from fixed;
    \xmark{} = does not. The Part 1 pattern baseline is not \tool; see
    Section~\ref{sec:baselines}.}
  \label{tab:rq2-comparison}
  \begin{tabular}{@{}lccc@{}}
    \toprule
    Tool & \shortstack{2022\\4450} & \shortstack{2023\\0215} & \shortstack{2024\\4741} \\
    \midrule
    Part 1 pattern baseline & \cmark & \cmark & \xmark$^{a}$ \\
    Clang Static Analyzer~18 & \xmark & \xmark & \xmark \\
    Coccinelle~1.1.1 & \xmark & \xmark & \xmark \\
    Cppcheck~2.13.0 & \xmark & \xmark & \xmark \\
    \bottomrule
  \end{tabular}
\end{table}

The columns correspond to CVE-2022-4450~\cite{nvd2022cve4450},
CVE-2023-0215~\cite{nvd2023cve0215}, and
CVE-2024-4741~\cite{nvd2024cve4741}, respectively. $^{a}$On the
corrected CVE-2024-4741 pair (\texttt{ec559c2a6c}/\texttt{b3f0eb0a29}) the
Part~1 baseline alerts on both sides; it is quiet only on the later hardening
commit \texttt{2d05959073}, so it matches that pattern, not the advisory's
state-guard fix.
Table~\ref{tab:rq2-comparison} compares four tools on six pinned revisions
(three CVEs, vulnerable and fixed side). A Clang positive control, run on a
synthetic use-after-free, produced the expected diagnostic with the same
configuration, confirming the analyzer was active for the clean runs.

\begin{observation}
\label{obs:rq2-generic-miss}
None of three general-purpose open-source tools, run with their stock
rules and no CVE-specific tuning, distinguishes the vulnerable revision
from the fixed revision for any of the three mechanisms. Coccinelle's
generic release-event rule fires near \texttt{BIO\_new\_NDEF} in both
revisions, but does not inspect \texttt{BIO\_pop} or the linked-BIO state,
so it does not change between versions; Cppcheck's output on the three
files is unrelated style diagnostics, identical in count before and after
each fix; Clang emits no diagnostic in any of the six runs.
\end{observation}

This does not show \tool beats these tools: Protocol B (\tool itself) does
not implement a rule for two of these three mechanisms
(Section~\ref{sec:subjects}). It shows that a source check specific to one
ownership or reset convention is needed at all, because the general-purpose
tools' stock configuration does not substitute for it.

\subsection{\rqref{3}: Does the rule generalize to an unseen later fix?}
\label{sec:rq3}

The tree at \texttt{d2217c88df} still flags \texttt{pvk\_set\_membuf} in
\texttt{pvkkdf.c}. OpenSSL fixed it one day later, in \texttt{a632bfaa4e},
a direct child of \texttt{d2217c88df}, with the same \texttt{*buffer =
NULL} change as the other four KDF sites.

\begin{observation}
\label{obs:rq3-timetravel}
\tool, run unchanged from its CVE-2022-4450 derivation, flags a site that a
different, later commit by OpenSSL went on to fix. This is one site, not a
rate, and we report it as a single encouraging case, not as a time-travel
precision or recall figure: a systematic time-travel evaluation over more
sites is future work.
\end{observation}

\subsection{\rqref{4}: What does \tool cost in time?}
\label{sec:rq4}

\begin{observation}
\label{obs:rq4-cost}
The independent rerun (Section~\ref{sec:setup}) reproduced every count in
Table~\ref{tab:rq1-wholetree} exactly. Across five full-tree scans of
999--1{,}155 files plus two single-function checks, CPU time was under 48
seconds (27.17\,s user, 20.34\,s system) of 5~min 31~s wall-clock; the rest
is one-time network latency from a partial clone of OpenSSL. The CPU figure
still includes \texttt{git} sparse-checkout cost, which we did not
instrument separately.
\end{observation}

\subsection{\rqref{5}: Does a second, differently-shaped rule generalize the same way?}
\label{sec:rq5}

Protocol A checks a different invariant shape than Protocol B:
establish-then-required-detach-before-release on a local variable rather
than release-then-required-reset on a parameter, with a syntactic alias
rule for the common \texttt{x~=~BIO\_push(v,~\ldots{})} idiom
(Section~\ref{sec:subjects}). It reuses Protocol B's \texttt{tree-sitter}
parser and CFG; only the event rule and path search are new. We ran it
two ways, both independently reproduced on a separate machine
(Section~\ref{sec:setup}).

\begin{observation}
\label{obs:rq5-protocola-motivating}
On \texttt{BIO\_new\_NDEF}, Protocol A flags \texttt{asn\_bio} as a
violation before the CVE-2023-0215 fix and as conforming after it, the
same pre/post-fix distinction Protocol B makes for CVE-2022-4450. The
rule-as-data design extends to a second invariant shape without changing
the parser or the CFG builder.
\end{observation}

On the same master snapshot, Protocol~A flags 11 sites; we checked every
one against the pinned source (\path{results/protocol_a_triage.md}).

\begin{observation}
\label{obs:rq5-protocola-master}
0 of the 11 whole-tree violations are confirmed or plausible recurrences.
One is a tool limitation: a real \texttt{BIO\_pop} is guarded by a boolean
flag the path search cannot correlate with the branch that set it. The
other 10 release a chain built entirely from freshly-allocated local
objects (two behind a push whose failure branch is unreachable once its
first argument is proven non-\texttt{NULL}). None splices in an
externally-owned BIO -- the shape that makes CVE-2023-0215 dangerous -- and
the rule has no notion of ownership to tell them apart.
\end{observation}

This confirms the gap named in Section~\ref{sec:related}: without ownership
or escape tracking, a single-object rule cannot generalize a BIO-chain
invariant across a codebase, even though the same architecture generalized
a release/reset invariant in \rqref{1}.
